%% physics_guide.tex — multiSpeciesRegionFoam Physics Reference
%% Compile:  pdflatex physics_guide.tex  (run twice for cross-references)
%%           or:  latexmk -pdf physics_guide.tex

\documentclass[11pt,a4paper]{article}

%% Packages
\usepackage[top=25mm,bottom=25mm,left=25mm,right=25mm]{geometry}
\usepackage{amsmath,amssymb,mathtools}
\usepackage{booktabs,array,tabularx}
\usepackage{graphicx}
\usepackage{xcolor}
\usepackage{tcolorbox}
\tcbuselibrary{skins,breakable}
\usepackage{hyperref}
\usepackage{cleveref}
\usepackage{enumitem}
\usepackage{microtype}
\usepackage{fancyhdr}
\usepackage{titlesec}
\usepackage{parskip}
\usepackage{siunitx}
\usepackage[T1]{fontenc}
\usepackage{lmodern}
\usepackage{multirow}
\usepackage{float}
\usepackage{longtable}

%% Colours
\definecolor{foamblue}{RGB}{0,82,147}
\definecolor{casegreen}{RGB}{30,120,60}
\definecolor{eqngray}{RGB}{245,245,248}
\definecolor{warnorange}{RGB}{200,100,0}

%% Hyperref
\hypersetup{
    colorlinks=true,
    linkcolor=foamblue,
    citecolor=foamblue,
    urlcolor=foamblue,
    pdftitle={multiSpeciesRegionFoam — Physics Reference Guide},
    pdfauthor={Jaydeep Deshpande},
}

%% Section formatting
\titleformat{\section}{\large\bfseries\color{foamblue}}{}{0em}{}[\color{foamblue}\titlerule]
\titleformat{\subsection}{\normalsize\bfseries\color{foamblue!80}}{}{0em}{}
\titleformat{\subsubsection}{\normalsize\itshape}{}{0em}{}

%% tcolorbox styles
\tcbset{
  eqbox/.style={
    enhanced, breakable,
    colback=eqngray, colframe=foamblue!60,
    arc=2pt, boxrule=0.6pt,
    left=6pt, right=6pt, top=4pt, bottom=4pt,
    fonttitle=\bfseries\small\color{foamblue},
    attach boxed title to top left={yshift=-2mm,xshift=4mm},
    boxed title style={colback=foamblue!15,colframe=foamblue!60,arc=2pt,boxrule=0.4pt},
  },
  casebox/.style={
    enhanced, breakable,
    colback=casegreen!5, colframe=casegreen!60,
    arc=2pt, boxrule=0.6pt,
    left=6pt, right=6pt, top=4pt, bottom=4pt,
  },
  notebox/.style={
    enhanced,
    colback=warnorange!8, colframe=warnorange!60,
    arc=2pt, boxrule=0.6pt,
    left=6pt, right=6pt, top=4pt, bottom=4pt,
  }
}

%% Header / footer
\pagestyle{fancy}
\fancyhf{}
\fancyhead[L]{\small\color{foamblue}\textit{multiSpeciesRegionFoam — Physics Reference}}
\fancyhead[R]{\small\color{foamblue}\thepage}
\fancyfoot[C]{\small\color{gray}OpenFOAM-13 library for multi-region species transport}
\renewcommand{\headrulewidth}{0.4pt}

%% Notation shortcuts
\newcommand{\pder}[2]{\frac{\partial #1}{\partial #2}}
\newcommand{\mol}{\,\text{mol}}
\newcommand{\m}{\,\text{m}}
\newcommand{\K}{\,\text{K}}
\newcommand{\s}{\,\text{s}}
\newcommand{\Pa}{\,\text{Pa}}
\newcommand{\J}{\,\text{J}}
\newcommand{\W}{\,\text{W}}
\newcommand{\kg}{\,\text{kg}}
\newcommand{\caseref}[1]{\texttt{\color{casegreen}#1}}
\newcommand{\bctype}[1]{\texttt{\color{foamblue}#1}}
\newcommand{\Da}{\mathrm{Da}}

% ============================================================
\begin{document}
% ============================================================

\begin{titlepage}
  \centering
  \vspace*{20mm}
  {\color{foamblue}\rule{\linewidth}{2pt}}\\[6mm]
  {\Huge\bfseries\color{foamblue} multiSpeciesRegionFoam}\\[3mm]
  {\Large Physics Reference Guide}\\[4mm]
  {\color{foamblue}\rule{\linewidth}{1pt}}\\[10mm]
  {\large Species transport through solid membranes and porous materials\\
   in the OpenFOAM-13 \texttt{foamMultiRun} framework}\\[15mm]

  \begin{tcolorbox}[casebox,width=0.85\linewidth]
    \textbf{Scope.}
    This document describes every physical model implemented in the
    \texttt{multiSpeciesRegionFoam} library: governing equations, material
    constitutive laws, interface and surface boundary conditions, performance
    metrics, and numerical discretisation. Each section indicates which tutorial
    cases (case311–case330) demonstrate the corresponding physics.
  \end{tcolorbox}

  \vspace{30pt}


  {\large\bfseries Jaydeep Deshpande}\\[6mm]
  \vfill
  {\small OpenFOAM-13 $\cdot$ \href{https://openfoam.org}{openfoam.org}}
\end{titlepage}

\tableofcontents
\newpage

% ============================================================
\section{Introduction and Scope}
\label{sec:intro}
% ============================================================

\texttt{multiSpeciesRegionFoam} is an OpenFOAM-13 add-on that extends the
native \texttt{foamMultiRun}~\cite{CFDDirect} multi-region framework with coupled species
transport.  Two solver modules are provided:

\begin{itemize}[itemsep=2pt]
  \item \textbf{\texttt{speciesSolid}} — adds a species conservation equation to solid
        (or fluid-region-with-prescribed-flow) regions alongside the inherited thermal
        energy equation from the \texttt{solid} base module.
  \item \textbf{\texttt{speciesFluid}} — extends \texttt{incompressibleFluid} with a
        solved temperature field and species transport equation.  The full
        Navier--Stokes PIMPLE loop (U, p) is inherited; temperature and species
        are added in \texttt{thermophysicalPredictor()}.
  \item \textbf{\texttt{compressibleSpeciesFluid}} — extends the OpenFOAM-13
        \texttt{fluid} module (compressible PIMPLE with full internal-energy
        equation) with species transport.  Required when density varies with
        pressure, or when the standard OpenFOAM \texttt{heRhoThermo} property
        infrastructure is preferred (e.g.\ for gases or using \texttt{rhoConst}
        for liquids with thermophysical transport).
\end{itemize}

Inter-region coupling (heat and species) uses OpenFOAM's mapped-patch
infrastructure, and all new boundary conditions are registered in the
standard run-time selection tables~\cite{Hattab2024}.

\textbf{Physical scope.}
The library covers the full range from pure-diffusion permeation barriers to
coupled advection-diffusion in molten-salt channels.  Typical applications are:

\begin{itemize}[itemsep=2pt]
  \item Hydrogen / tritium permeation through first-wall and breeding-blanket
        materials in fusion machines; structural materials in molten salt
        reactors (Sieverts-law solubility).
  \item Dense-film membrane separation: reverse osmosis, membrane distillation,
        pervaporation (Henry-law solubility, Arrhenius temperature dependence).
  \item Permeation barriers (tungsten carbide, alumina coatings)~\cite{Anderl1992}.
  \item Heat-exchanger tritium transport: FLiBe/Hitec molten-salt systems with
        conjugate heat transfer and dissolved-species permeation.
  \item Fluid-channel species accumulation and wall permeation (speciesFluid).
\end{itemize}

% ============================================================
\section{Governing Equations}
\label{sec:gov}
% ============================================================

\subsection{Species conservation}
\label{sec:species}

In each solid region the \emph{mobile} species concentration
$C\,[\mol\,\m^{-3}]$ satisfies:

\begin{tcolorbox}[eqbox,title={Species equation — every solid region}]
\begin{equation}\label{eq:species}
  \pder{C}{t}
  \;=\;
  \nabla\cdot\!\bigl(D(T)\,\nabla C\bigr)
  \;-\; \pder{C_t}{t}
  \;+\; S_\text{vol}
\end{equation}
\end{tcolorbox}

\begin{center}
\begin{tabular}{lll}
\toprule
Symbol & Meaning & Units \\
\midrule
$C$       & Mobile species concentration      & $\mol\,\m^{-3}$ \\
$D(T)$    & Temperature-dependent diffusivity  & $\m^2\,\s^{-1}$ \\
$C_t$     & Trapped concentration (§\ref{sec:trap}) & $\mol\,\m^{-3}$ \\
$S_\text{vol}$ & Uniform volumetric source     & $\mol\,\m^{-3}\,\s^{-1}$ \\
\bottomrule
\end{tabular}
\end{center}

When trapping is disabled (\texttt{trappingModel noTrapping}), $\partial C_t/\partial t = 0$
and the equation reduces to pure diffusion.  The source term $S_\text{vol}$ is
optional and uniform within each region.

\begin{tcolorbox}[casebox]
  \textbf{Demonstrated by:} all twenty tutorial cases
  (\caseref{case311}–\caseref{case330}).
  Cases \caseref{case311}–\caseref{case312} use the thermal field $T$ as a
  species proxy to verify the pure-diffusion solver before species coupling is
  introduced.
\end{tcolorbox}

\subsection{Energy conservation}
\label{sec:energy}

The thermal energy equation is inherited from OpenFOAM-13's \texttt{solid}
module:

\begin{tcolorbox}[eqbox,title={Energy equation (Fourier conduction)}]
\begin{equation}\label{eq:energy}
  \pder{(\rho\, e)}{t}
  \;+\; \nabla\cdot\mathbf{q}
  \;=\; S_e,
  \qquad
  \mathbf{q} = -\kappa(T)\,\nabla T
\end{equation}
\end{tcolorbox}

\begin{center}
\begin{tabular}{lll}
\toprule
Symbol & Meaning & Units \\
\midrule
$\rho$     & Density                     & $\text{kg}\,\m^{-3}$ \\
$e$        & Specific internal energy     & $\J\,\text{kg}^{-1}$ \\
$\mathbf{q}$ & Conductive heat flux       & $\W\,\m^{-2}$ \\
$\kappa(T)$ & Thermal conductivity        & $\W\,\m^{-1}\,\K^{-1}$ \\
$S_e$      & Volumetric heat source       & $\W\,\m^{-3}$ \\
\bottomrule
\end{tabular}
\end{center}

The species and energy equations are solved \emph{sequentially} within each
PIMPLE outer corrector.  After the energy equation updates $T$, the
diffusivity $D(T)$ (and all Arrhenius properties) are refreshed before
assembling the species equation.

\begin{tcolorbox}[casebox]
  \textbf{Full energy–species coupling:}
  \caseref{case316} (membrane distillation with Arrhenius $D(T)$ driven by a
  temperature gradient), \caseref{case320} (three-region shell-and-tube
  heat exchanger; conjugate heat transfer through FLiBe–SS316–Hitec).\\
  \textbf{Isothermal (energy solved but $T$ fixed):}
  \caseref{case311}–\caseref{case315}, \caseref{case317}–\caseref{case319},
  \caseref{case321}–\caseref{case322}, \caseref{case328}.\\
  \textbf{Active $T$ advection (\texttt{speciesFluid}):}
  \caseref{case327}–\caseref{case328}.\\
  \textbf{Compressible energy eq.\ (\texttt{compressibleSpeciesFluid}):}
  \caseref{case329}–\caseref{case330}.
\end{tcolorbox}

\subsection{Species and energy transport in fluid regions (\texttt{speciesFluid})}
\label{sec:fluid}

The \texttt{speciesFluid} solver module targets incompressible fluid regions
where the velocity field must be computed from the Navier--Stokes equations.
It inherits the PIMPLE velocity-pressure loop from \texttt{incompressibleFluid}
and adds temperature and species transport via an overridden
\texttt{thermophysicalPredictor()} method.

\subsubsection{Full equation set solved per PIMPLE outer corrector}

\textbf{Step 1 — Momentum and continuity} (inherited, solved first):

\begin{tcolorbox}[eqbox,title={Incompressible Navier--Stokes}]
\begin{align}
  \nabla\cdot\mathbf{U} &= 0 \label{eq:cont} \\[4pt]
  \pder{\mathbf{U}}{t}
    + \nabla\cdot(\mathbf{U}\mathbf{U})
  &= -\nabla p + \nabla\cdot(\nu\,\nabla\mathbf{U}) \label{eq:mom}
\end{align}
\end{tcolorbox}

where $p\,[\m^2\,\s^{-2}]$ is kinematic pressure and $\nu\,[\m^2\,\s^{-1}]$
is kinematic viscosity.

\textbf{Step 2 — Temperature} (scalar transport):

\begin{tcolorbox}[eqbox,title={Energy equation — \texttt{speciesFluid} regions}]
\begin{equation}\label{eq:fluid_T}
  \pder{T}{t}
  \;+\; \nabla\cdot(\mathbf{U}\,T)
  \;=\; \nabla\cdot\!\bigl(\alpha\,\nabla T\bigr),
  \qquad
  \alpha = \frac{\kappa}{\rho\,C_p}
\end{equation}
\end{tcolorbox}

This is the incompressible constant-property form of the energy equation,
where $\alpha\,[\m^2\,\s^{-1}]$ is the thermal diffusivity.  It is solved as
an advection-diffusion equation for $T$ directly, rather than for the internal
energy $e = C_p\,T$ used by the \texttt{solid} module.

\textbf{Step 3 — Species concentration} (advection-diffusion):

\begin{tcolorbox}[eqbox,title={Species equation — \texttt{speciesFluid} regions}]
\begin{equation}\label{eq:fluid_C}
  \pder{C}{t}
  \;+\; \nabla\cdot(\mathbf{U}\,C)
  \;=\; \nabla\cdot\!\bigl(D(T)\,\nabla C\bigr)
  \;-\; \pder{C_t}{t}
  \;+\; S_\text{vol}
\end{equation}
\end{tcolorbox}

After solving $T$, the diffusivity $D(T)$ is refreshed via
\texttt{species\_.correctProperties()} before assembling~\eqref{eq:fluid_C}.
The trapping sink $\partial C_t/\partial t$ (§\ref{sec:trap}) applies equally
here; most fluid-region cases use \texttt{noTrapping}.

\subsubsection{Numerical requirements for \texttt{speciesFluid}}

The advection terms $\nabla\cdot(\mathbf{U}\,T)$ and $\nabla\cdot(\mathbf{U}\,C)$
produce \emph{asymmetric} coefficient matrices.  The following requirements
must be observed in \texttt{fvSolution}:

\begin{tcolorbox}[notebox]
\textbf{Matrix solver:} use \texttt{PBiCGStab} with \texttt{DILU} preconditioner
for \texttt{"T.*"} and \texttt{"C\_.*"}.  \texttt{PCG} (symmetric-only) will
abort with the error \textit{``Unknown asymmetric matrix solver PCG''}.

\textbf{Wildcard:} use \texttt{"T.*"} (not bare \texttt{T}) in the
\texttt{solvers} block.  The PIMPLE final corrector creates a \texttt{TFinal}
field; omitting the wildcard causes a \textit{``keyword TFinal is undefined''} error.
\end{tcolorbox}

\textbf{Minimum \texttt{fvSolution} for a \texttt{speciesFluid} region:}
\begin{verbatim}
solvers {
    "T.*"  { solver PBiCGStab; preconditioner DILU;
              tolerance 1e-10; relTol 0; }
    "C_.*" { solver PBiCGStab; preconditioner DILU;
              tolerance 1e-10; relTol 0; }
}
PIMPLE { nOuterCorrectors 3; nCorrectors 2;
         nNonOrthogonalCorrectors 0; }
\end{verbatim}

\textbf{Required \texttt{physicalProperties} format:}
\begin{verbatim}
viscosityModel  Newtonian;           % read by incompressibleFluid
nu  [0 2 -1 0 0 0 0]  4.74e-7;
rho    983;                          % read by speciesFluid
Cp     4183;
mixture { transport { kappa  0.654; } }
\end{verbatim}

\begin{tcolorbox}[casebox]
  \textbf{Demonstrated by (\texttt{speciesFluid}):}
  \caseref{case327} (DCMD with solved laminar flow; feed and permeate use
  \texttt{speciesFluid}),
  \caseref{case328} (Poiseuille + Graetz validation of the solver).\\
  \textbf{See also:} \caseref{case329}–\caseref{case330} use
  \texttt{compressibleSpeciesFluid} (§\ref{sec:compfluid}), which wraps the
  same PIMPLE loop via the compressible \texttt{fluid} base class.
\end{tcolorbox}

\subsubsection{Dimensional considerations — Schmidt and Péclet numbers}

In fluid regions, two dimensionless groups govern the relative importance of
advection and diffusion for momentum ($\mathrm{Sc}$) and species
($\mathrm{Pe}$):

\begin{equation}
  \mathrm{Sc} = \frac{\nu}{D}, \qquad
  \mathrm{Pe} = \frac{U\,L}{D} = \mathrm{Re}\cdot\mathrm{Sc}
\end{equation}

For liquid metals and molten salts, $\mathrm{Sc}\sim 10^2$--$10^3$ (molecular
diffusivity is orders of magnitude smaller than momentum diffusivity).  At high
$\mathrm{Pe}$ the concentration boundary layer at a permeation wall is thin:

\begin{equation}
  \frac{\delta_C}{L} \;\sim\; \mathrm{Pe}^{-1/3}
\end{equation}

This sets the mesh resolution requirement near permeation walls.
Upwind (\texttt{Gauss upwind}) or linearUpwind schemes for
$\nabla\cdot(\mathbf{U}\,C)$ prevent numerical oscillations at high Pe.

\subsection{Species transport in compressible fluid regions (\texttt{compressibleSpeciesFluid})}
\label{sec:compfluid}

The \texttt{compressibleSpeciesFluid} solver extends the OpenFOAM-13
\texttt{fluid} base module — which provides a compressible PIMPLE loop and
full internal-energy equation — with single-species molar concentration
transport.

\subsubsection{Full equation set per PIMPLE outer corrector}

\textbf{Step 1 — Continuity + momentum + pressure} (inherited from
\texttt{isothermalFluid}/\texttt{fluid}):

\begin{tcolorbox}[eqbox,title={Compressible continuity and momentum}]
\begin{align}
  \pder{\rho}{t} + \nabla\cdot(\rho\mathbf{U}) &= 0 \label{eq:cont_comp} \\[4pt]
  \pder{(\rho\mathbf{U})}{t} + \nabla\cdot(\rho\mathbf{U}\mathbf{U})
  &= -\nabla p + \nabla\cdot\boldsymbol{\tau} \label{eq:mom_comp}
\end{align}
\end{tcolorbox}

\textbf{Step 2 — Internal energy} (inherited, solved for $e$):

\begin{tcolorbox}[eqbox,title={Compressible energy equation (\texttt{fluid} base)}]
\begin{equation}\label{eq:energy_comp}
  \pder{(\rho e)}{t}
  + \nabla\cdot(\varphi\, e)
  + \nabla\cdot(\varphi\, K)
  + \nabla\cdot\!\left(\varphi\,\frac{p}{\rho}\right)
  + \nabla\cdot\mathbf{q}
  = S_e
\end{equation}
\end{tcolorbox}

Here $\varphi = \rho\mathbf{U}\cdot\mathbf{S}_f\,[\text{kg\,s}^{-1}]$ is the
face mass flux, $K = \tfrac{1}{2}|\mathbf{U}|^2$ is specific kinetic energy,
$\mathbf{q} = -\kappa\nabla T$ is the Fourier heat flux, and
$e\,[\J\,\text{kg}^{-1}]$ is the sensible internal energy.  After solving,
the temperature is recovered from the equation of state via
$T = T_\text{ref} + e/C_v$ (for \texttt{eConst}) and density is updated.

\textbf{Step 3 — Species concentration}:

\begin{tcolorbox}[eqbox,title={Species equation — \texttt{compressibleSpeciesFluid}}]
\begin{equation}\label{eq:comp_C}
  \pder{C}{t}
  + \nabla\cdot(\mathbf{U}_\text{vol}\,C)
  = \nabla\cdot\!\bigl(D(T)\,\nabla C\bigr)
  - \pder{C_t}{t}
  + S_\text{vol}
\end{equation}
\end{tcolorbox}

The volumetric flux $\mathbf{U}_\text{vol}$ is derived from the mass flux:
\begin{equation}\label{eq:phiU}
  \varphi_\text{vol} = \frac{\varphi}{\rho_f}
  = \frac{\phi}{\texttt{fvc::interpolate(rho)}}
  \quad [\text{m}^3\,\text{s}^{-1}]
\end{equation}

\begin{tcolorbox}[notebox]
\textbf{Critical:} the mass flux $\varphi\,[\text{kg\,s}^{-1}]$ cannot be used
directly in the species equation because $C$ is molar concentration
$[\mol\,\m^{-3}]$, not mass concentration.  Using $\varphi$ would introduce a
spurious factor of $\rho$ in the advection term.  The conversion
$\varphi_\text{vol} = \varphi/\rho_f$ is mandatory.
\end{tcolorbox}

\subsubsection{Thermophysical property format (\texttt{heRhoThermo})}

Unlike \texttt{speciesFluid}, which reads a flat \texttt{physicalProperties}
dictionary, \texttt{compressibleSpeciesFluid} requires the full OpenFOAM
\texttt{thermoType} hierarchy in \texttt{constant/<region>/physicalProperties}:

\begin{verbatim}
thermoType {
    type            heRhoThermo;
    mixture         pureMixture;
    transport       const;
    thermo          eConst;
    equationOfState rhoConst;   % or perfectGas for true compressibility
    specie          specie;
    energy          sensibleInternalEnergy;
}
mixture {
    specie          { molWeight  <M [g/mol]>; }
    equationOfState { rho  <rho [kg/m3]>; }    % for rhoConst
    thermodynamics  { Cv  <Cv [J/kg/K]>;  Hf  0; }
    transport       { mu  <mu [Pa.s]>;  Pr  <Pr>; }
}
\end{verbatim}

A second dictionary \texttt{constant/<region>/thermophysicalTransport} is also
required:
\begin{verbatim}
laminar { model Fourier; }
\end{verbatim}

\subsubsection{Required \texttt{fvSchemes} and \texttt{fvSolution}}

\textbf{fvSchemes} — three additional \texttt{divSchemes} entries compared to
\texttt{speciesFluid}:
\begin{verbatim}
divSchemes {
    div(phi,e)          Gauss linearUpwind grad(e);
    div(phi,K)          Gauss linear;
    div(phi,(p|rho))    Gauss linear;
    div(phiU,C_H2)      Gauss upwind;   % phiU = phi/rho_face
}
\end{verbatim}

\textbf{fvSolution} — energy field is \texttt{e} (not \texttt{T}); density is
updated explicitly via the EOS:
\begin{verbatim}
solvers {
    "e.*"   { solver PBiCGStab; preconditioner DILU;
              tolerance 1e-10; relTol 0; }
    "C_.*"  { solver PBiCGStab; preconditioner DILU;
              tolerance 1e-10; relTol 0; }
    "rho.*" { solver diagonal; }      % covers rhoFinal on final corrector
}
\end{verbatim}

Pressure uses thermodynamic dimensions $[\text{Pa}]$ (not kinematic
$[\m^2\,\s^{-2}]$): \texttt{dimensions [1 -1 -2 0 0 0 0]}, with
\texttt{fixedValue} at the outlet (e.g.\ 101325\,Pa for atmospheric reference).

\begin{tcolorbox}[casebox]
  \textbf{Demonstrated by:}
  \caseref{case329} (FLiBe molten-salt channel at 873\,K with distributed
  H$_2$ source and SS316 permeation; \texttt{rhoConst} EOS validates
  equivalence to the incompressible \texttt{speciesFluid}),
  \caseref{case330} (He gas channel at 700\,K, 1\,atm; same source, $D$,
  and $K_s$ as \caseref{case329} for cross-solver flux comparison).
\end{tcolorbox}

\subsection{Species and energy coupling — multi-region}
\label{sec:coupling}

In multi-region cases, the PIMPLE loop exchanges temperature across coupled
patches via \bctype{coupledTemperature} boundary conditions (standard
OpenFOAM), and exchanges species concentration via the custom
\bctype{sievertsCoupledMixed} or \bctype{surfaceRecombination} BCs described
in §\ref{sec:bc}.  Both exchanges happen every PIMPLE outer corrector, so the
fields remain thermodynamically consistent throughout the time step.

% ============================================================
\section{Material Models}
\label{sec:mat}
% ============================================================

\subsection{Arrhenius temperature dependence}
\label{sec:arrhenius}

Every material property that depends on temperature follows the Arrhenius form.
This applies uniformly to $D$, the solubility $K_s$, and the surface-kinetics
coefficients $K_d$, $K_r$:

\begin{tcolorbox}[eqbox,title={Arrhenius property}]
\begin{equation}\label{eq:arrhenius}
  X(T) \;=\; X_0 \;\exp\!\left(-\frac{E_a}{R\,T}\right)
\end{equation}
\end{tcolorbox}

\begin{center}
\begin{tabular}{lll}
\toprule
Symbol & Meaning & Notes \\
\midrule
$X_0$  & Pre-exponential factor & Carries the SI units of $X$ \\
$E_a$  & Activation energy      & $\J\,\mol^{-1}$ \\
$R$    & Gas constant           & $8.314\,\J\,\mol^{-1}\,\K^{-1}$ \\
$T$    & Local temperature      & $\K$ \\
\bottomrule
\end{tabular}
\end{center}

Setting $E_a = 0$ gives a temperature-independent constant, used in all
isothermal tutorial cases.

\textbf{Dictionary syntax} (applies to all Arrhenius properties):
\begin{verbatim}
    D { X0  [0 2 -1 0 0 0 0]  2.9e-7;   Ea  13800; }
\end{verbatim}

\begin{tcolorbox}[casebox]
  \textbf{Isothermal ($E_a = 0$):} \caseref{case311}–\caseref{case315},
  \caseref{case317}–\caseref{case319}, \caseref{case321}–\caseref{case322},
  \caseref{case328}–\caseref{case330}.\\
  \textbf{Arrhenius $D(T)$ active:} \caseref{case316} (vapour diffusion),
  \caseref{case319} (WC and SS316), \caseref{case320} (SS316 in HX),
  \caseref{case323}–\caseref{case327} (PTFE membrane and/or SS316 wall).
\end{tcolorbox}

\subsection{Solubility laws}
\label{sec:sol}

\subsubsection{Sieverts' law (metals and interstitial solids)}

Sieverts' law describes the equilibrium solubility of diatomic gases (H$_2$,
D$_2$, T$_2$) that dissociate upon dissolution.  The dissolved concentration
is proportional to the \emph{square root} of the gas partial pressure:

\begin{tcolorbox}[eqbox,title={Sieverts' law}]
\begin{equation}\label{eq:sieverts}
  C \;=\; K_s(T)\;\sqrt{p}
\end{equation}
\end{tcolorbox}

$K_s(T)\,[\mol\,\m^{-3}\,\Pa^{-1/2}]$ is the Sieverts solubility constant,
following the Arrhenius form~\eqref{eq:arrhenius}.  Applicable materials
include iron, stainless steel (SS316), tungsten, vanadium alloys, and palladium~\cite{Ockwig2007}.

\begin{tcolorbox}[casebox]
  Used in: \caseref{case312}–\caseref{case314}, \caseref{case318}–\caseref{case320},
  \caseref{case321} (metal layer), \caseref{case322},
  \caseref{case323} (FLiBe and SS316 walls),
  \caseref{case329}–\caseref{case330} (fluid–SS316 Sieverts interface).
\end{tcolorbox}

\subsubsection{Henry's law (polymers, liquids, porous media)}

For species that dissolve without dissociation (vapour in a polymer, dissolved
gas in a liquid), the equilibrium concentration is proportional to pressure:

\begin{tcolorbox}[eqbox,title={Henry's law}]
\begin{equation}\label{eq:henry}
  C \;=\; K_H(T)\;p
\end{equation}
\end{tcolorbox}

$K_H(T)\,[\mol\,\m^{-3}\,\Pa^{-1}]$ is the Henry constant.  Applicable to
dense polymer films, membrane distillation membranes, and organic solvents~\cite{Baker2012}.

\begin{tcolorbox}[casebox]
  Used in: \caseref{case315}–\caseref{case316} (membrane distillation),
  \caseref{case317} (reverse osmosis, $K_H = 0.5$ at the feed/membrane
  interface), \caseref{case321} (polymer layer).
\end{tcolorbox}

\subsection{McNabb--Foster trapping kinetics}
\label{sec:trap}

In metals, hydrogen isotopes can occupy lattice defects (vacancies,
dislocations, grain boundaries).  The McNabb--Foster model~\cite{McNabb1963}
treats a single population of trap sites with density
$n_t = f_t N\,[\mol\,\m^{-3}]$:

\begin{tcolorbox}[eqbox,title={McNabb--Foster ODE}]
\begin{equation}\label{eq:trap}
  \pder{C_t}{t}
  \;=\;
  \frac{\alpha_t}{N}\,C\,(n_t - C_t)
  \;-\;
  \alpha_d(T)\,C_t
\end{equation}
\end{tcolorbox}

\begin{center}
\begin{tabular}{lll}
\toprule
Symbol & Meaning & Units \\
\midrule
$N$      & Molar trap-site density               & $\mol\,\m^{-3}$ \\
$f_t$    & Trap fraction ($0 < f_t \le 1$)        & --- \\
$n_t = f_t N$ & Available trap-site density      & $\mol\,\m^{-3}$ \\
$\alpha_t$ & Trapping rate coefficient             & $\s^{-1}$ \\
$\alpha_d(T) = \alpha_{d0}\exp(-\varepsilon/k_B T)$ & Detrapping rate & $\s^{-1}$ \\
$\varepsilon$ & Trap binding energy               & eV \\
$k_B$    & Boltzmann constant                     & $\J\,\K^{-1}$ \\
\bottomrule
\end{tabular}
\end{center}

The trapping term $\partial C_t/\partial t$ is subtracted from the mobile
species equation~\eqref{eq:species} as a sink.

\subsubsection{Trapping regimes}

Two limiting regimes exist depending on the ratio of trapping to diffusion time scales:

\begin{itemize}[itemsep=2pt]
  \item \textbf{Weak trapping} ($\alpha_t / D \ll 1$): traps fill slowly; the
        mobile front propagates at nearly the bare diffusion speed; breakthrough
        time $\approx L^2/D$.
  \item \textbf{Strong trapping} ($\alpha_t / D \gg 1$): most atoms are trapped;
        effective diffusivity $D_\text{eff} = D/(1 + N_t\,\alpha_t)$ is greatly
        reduced; breakthrough is delayed by orders of magnitude.
\end{itemize}

\subsubsection{Numerical discretisation of the trapping ODE}

The ODE~\eqref{eq:trap} is stiff and must be discretised carefully to remain
bounded ($0 \le C_t \le n_t$).  A per-cell semi-implicit Euler step is used:

\begin{tcolorbox}[eqbox,title={Semi-implicit Euler for trapping ODE}]
\begin{equation}\label{eq:trap_euler}
  C_t^{n+1}
  \;=\;
  \frac{C_t^n \;+\; \Delta t\,\dfrac{\alpha_t}{N}\,C\,n_t}
       {1 \;+\; \Delta t\!\left(\dfrac{\alpha_t}{N}\,C \;+\; \alpha_d\right)}
\end{equation}
\end{tcolorbox}

This scheme is unconditionally stable and guarantees $0 \le C_t^{n+1} \le n_t$
for any time step.

\begin{tcolorbox}[casebox]
  \textbf{Demonstrated by:} \caseref{case313}.
  Two sub-cases span the weak and strong trapping regimes.  The analytical
  breakthrough time is used to validate both limits.
\end{tcolorbox}

% ============================================================
\section{Interface and Surface Boundary Conditions}
\label{sec:bc}
% ============================================================

\subsection{Overview}

At every coupled patch, two physical conditions must hold simultaneously:

\begin{enumerate}[itemsep=2pt]
  \item \textbf{Flux continuity:}
        $D_A\,(\partial C_A/\partial n)_\text{face} =
         D_B\,(\partial C_B/\partial n)_\text{face}$
  \item \textbf{Thermodynamic equilibrium:} determined by the solubility
        law on each side (§\ref{sec:sol}).
\end{enumerate}

These are implemented as a \emph{mixed} (Robin) boundary condition in
OpenFOAM, derived from the conjugate heat-transfer analogy.

\subsection{Sieverts--Sieverts interface (linear partition)}
\label{sec:bc_linear}

When both sides follow Sieverts' law, the equilibrium condition at the
interface face is:

\begin{tcolorbox}[eqbox,title={Sieverts--Sieverts partition (\texttt{partition linear})}]
\begin{equation}\label{eq:linear_part}
  \frac{C_A}{K_{s,A}} \;=\; \frac{C_B}{K_{s,B}}
  \qquad\Longleftrightarrow\qquad
  C_A \;=\; \frac{K_{s,A}}{K_{s,B}}\,C_B
\end{equation}
\end{tcolorbox}

The concentration ratio $K_{s,A}/K_{s,B}$ can be greater or less than unity,
producing a concentration jump at the interface even in steady state.  When
$K_{s,A} = K_{s,B}$, the concentration is continuous.

\textbf{Mixed-BC coefficients} (self = region $A$, neighbour = region $B$):
\begin{align}
  \text{refValue} &= \frac{K_{s,A}}{K_{s,B}}\,C_{B,\text{cell}} \notag\\
  w &= \frac{D_B\,\delta_B^{-1}}{D_B\,\delta_B^{-1} + D_A\,\delta_A^{-1}(K_{s,A}/K_{s,B})} \label{eq:weight_linear}\\
  C_{A,\text{face}} &= w\cdot\text{refValue} + (1-w)\,C_{A,\text{cell}} \notag
\end{align}
where $\delta^{-1}$ is the distance from the cell centre to the interface face.

\begin{tcolorbox}[casebox]
  \textbf{Demonstrated by:} \caseref{case314} (two-layer composite membrane,
  $K_s$ ratio $= 4$; interface concentration jump of factor 4 at steady state),
  \caseref{case318} (code-to-code benchmark, equal $K_s$, continuous profile),
  \caseref{case319} (WC/SS316 barrier, large $K_s$ ratio responsible for PRF),
  \caseref{case320} (SS316 wall; also \caseref{case321} reference sub-case).
\end{tcolorbox}

\subsection{Henry--Sieverts interface (nonlinear partition)}
\label{sec:bc_nonlinear}

When a Henry-law material (polymer, liquid) interfaces with a Sieverts-law
material (metal), the same partial pressure $p_\text{int}$ must hold on both
sides:
\[
  p_\text{int} = \frac{C_\text{Henry}}{K_H} = \left(\frac{C_\text{Sieverts}}{K_s}\right)^{\!2}
\]

This gives two equivalent partition conditions depending on which side is
considered ``self'':

\begin{tcolorbox}[eqbox,title={Henry--Sieverts partition}]
\begin{align}
  \label{eq:sqrt_part}
  \text{Sieverts side (\texttt{partition sqrt}):} \quad
  &C_\text{Sieverts} = K_s \sqrt{\dfrac{C_\text{Henry}}{K_H}}\\[4pt]
  \label{eq:quad_part}
  \text{Henry side (\texttt{partition quadratic}):} \quad
  &C_\text{Henry} = K_H \left(\dfrac{C_\text{Sieverts}}{K_s}\right)^{\!2}
\end{align}
\end{tcolorbox}

Both sides must be set consistently in the \texttt{0/} boundary field files;
they encode the same physics from opposite perspectives.

\textbf{Linearised mixed-BC weight} for the nonlinear cases.
The weight~\eqref{eq:weight_linear} is derived by linearising the partition
function at the current Picard iterate.  For the \texttt{sqrt} mode
(self = Sieverts, neighbour = Henry):
\begin{equation}\label{eq:weight_sqrt}
  \frac{\mathrm{d}C_s}{\mathrm{d}C_n}\bigg|_\text{iterate}
  = \frac{C_{s,\text{face}}}{2\,C_{n,\text{cell}}}
  \qquad\Rightarrow\qquad
  \text{selfKD} = D_A\,\delta_A^{-1}\;\frac{C_{A,\text{face}}}{2\,C_{B,\text{cell}}}
\end{equation}

Using the linearised derivative (rather than the constant $K_s/K_H$) in the
weight formula ensures that flux continuity is enforced exactly at every
Picard step, even in steady state on a coarse mesh.

\subsubsection{Steady-state interface concentrations}

For the case of equal layer thicknesses $L$, equal diffusivities $D$, inlet
concentration $C_0 = 1\,\mol\,\m^{-3}$, and $K_s = K_H = 1$:

\begin{tcolorbox}[eqbox,title={Analytical interface concentrations (equal $D$, $L$, $K$)}]
\begin{equation}\label{eq:hs_ss}
  u^2 + u = 1,
  \quad u = C_{\text{Sieverts,int}},
  \qquad\Rightarrow\qquad
  u = \tfrac{\sqrt{5}-1}{2} \approx 0.618\,\mol\,\m^{-3}
\end{equation}
\end{tcolorbox}

The \emph{golden-ratio conjugate} $(\sqrt5-1)/2$ appears naturally because
the steady-state flux equality combined with the $\sqrt{\cdot}$ partition gives
exactly the quadratic $u^2+u-1=0$.  The Henry (polymer) side holds
$C_\text{Henry,int} = u^2 \approx 0.382\,\mol\,\m^{-3}$, demonstrating that the
\emph{concentration is discontinuous and higher on the Sieverts (metal) side}
— the defining signature of this partition law.

The steady-state flux in the Henry--Sieverts case is 23.6\,\% \emph{higher}
than the equivalent Sieverts--Sieverts case with the same $D$, $L$, and $K$:
\begin{equation}\label{eq:hs_flux}
  J_\text{HS} = D\,\frac{C_\text{Sieverts,int}}{L}
  = \frac{(\sqrt{5}-1)/2}{L/D}\,C_0
  \;\approx\; 1.236\;J_\text{SS}
\end{equation}

\begin{tcolorbox}[casebox]
  \textbf{Demonstrated by:} \caseref{case321}.
  Two sub-cases (\texttt{sieverts\_sieverts} reference and
  \texttt{henry\_sieverts}) run with identical $D$, $L$, boundary values;
  the 23.6\,\% flux difference is confirmed analytically and numerically.
\end{tcolorbox}

\subsection{Surface recombination and dissociation}
\label{sec:bc_recomb}

At a free metal surface exposed to a molecular gas, two kinetic processes
compete:

\begin{center}
\begin{tabular}{ll}
\textbf{Dissociative absorption:} & H$_2 \to 2$H,\quad $\text{flux} = K_d(T)\,p_\text{gas}$ \\
\textbf{Recombinative desorption:} & $2\text{H} \to \text{H}_2$,\quad $\text{flux} = K_r(T)\,C^2$ \\
\end{tabular}
\end{center}

The net flux into the solid is the \emph{surface recombination Robin BC}:

\begin{tcolorbox}[eqbox,title={Surface recombination BC (\texttt{surfaceRecombination})}]
\begin{equation}\label{eq:recomb}
  D\,\pder{C}{n}
  \;=\;
  K_d(T)\;p_\text{gas}
  \;-\;
  K_r(T)\;C^2
\end{equation}
\end{tcolorbox}

\begin{center}
\begin{tabular}{lll}
\toprule
Symbol & Meaning & Units \\
\midrule
$K_d(T)$ & Dissociation rate (Arrhenius) & $\mol\,\m^{-2}\,\s^{-1}\,\Pa^{-1}$ \\
$K_r(T)$ & Recombination rate (Arrhenius) & $\m^4\,\mol^{-1}\,\s^{-1}$ \\
$p_\text{gas}$ & Gas-phase partial pressure & $\Pa$ \\
\bottomrule
\end{tabular}
\end{center}

Both $K_d$ and $K_r$ follow the Arrhenius form~\eqref{eq:arrhenius}.
Setting $p_\text{gas} = 0$ models a vacuum or purge side (pure desorption).

\subsubsection{Equilibrium and Sieverts limit}

At equilibrium (zero net flux), $C = C_\text{eq}$:
\begin{equation}\label{eq:Ceq}
  K_d\,p_\text{gas} = K_r\,C_\text{eq}^2
  \qquad\Rightarrow\qquad
  C_\text{eq} = \sqrt{\frac{K_d\,p_\text{gas}}{K_r}}
\end{equation}

This is exactly Sieverts' law~\eqref{eq:sieverts} with $K_s = \sqrt{K_d/K_r}$:
\emph{the surface recombination BC generalises the Sieverts equilibrium BC and
recovers it in the fast-kinetics limit}~\cite{Pick1985}.

\subsubsection{Damköhler number}

The Damköhler number compares the surface-kinetics rate to the bulk-diffusion
rate:

\begin{tcolorbox}[eqbox,title={Damköhler number}]
\begin{equation}\label{eq:Da}
  \Da
  \;=\;
  \frac{K_r\,C_\text{eq}\,L}{D}
\end{equation}
\end{tcolorbox}

\begin{center}
\begin{tabular}{cll}
\toprule
$\Da$ & Limiting step & Surface concentration \\
\midrule
$\Da \gg 1$ & Diffusion-limited & $C_s \approx C_\text{eq}$ (Sieverts equilibrium) \\
$\Da \approx 1$ & Mixed control & $C_s = \tfrac{\sqrt{5}-1}{2}\,C_\text{eq}$ (see below) \\
$\Da \ll 1$ & Kinetically limited & $C_s \ll C_\text{eq}$ \\
\bottomrule
\end{tabular}
\end{center}

\subsubsection{Steady-state surface concentration}

At steady state, the Robin condition~\eqref{eq:recomb} must equal the
permeation flux $D\,C_s/L$ (assuming zero concentration at the purge face):
\begin{equation}\label{eq:recomb_ss}
  K_r\,C_s^2 + \frac{D}{L}\,C_s - K_d\,p_\text{gas} = 0
\end{equation}

The physical (positive) root is:
\begin{tcolorbox}[eqbox,title={Steady-state surface concentration (general Da)}]
\begin{equation}\label{eq:Cs_general}
  C_s
  \;=\;
  \frac{-D/L + \sqrt{(D/L)^2 + 4\,K_r\,K_d\,p_\text{gas}}}{2\,K_r}
\end{equation}
\end{tcolorbox}

\textbf{Special case $\Da = 1$} (with $K_d = K_r$, $p_\text{gas} = 1$,
$D/L = K_r\,C_\text{eq} = K_r$):
\begin{equation}\label{eq:Cs_Da1}
  C_s^2 + C_s - 1 = 0
  \qquad\Rightarrow\qquad
  C_s = \tfrac{\sqrt{5}-1}{2} \approx 0.618\,\mol\,\m^{-3}
\end{equation}

The golden-ratio conjugate appears again.  The flux is reduced by 38.2\,\%
relative to the Sieverts equilibrium limit:
\begin{equation}\label{eq:flux_Da1}
  \frac{J_\text{surf}}{J_\text{ref}} = \frac{C_s}{C_\text{eq}} = \frac{\sqrt{5}-1}{2}
  \approx 0.618
\end{equation}

\textbf{Numerical treatment.}
The $C^2$ term is Picard-frozen at the previous cell concentration,
making the gradient a simple field evaluation each PIMPLE outer iteration.
Convergence is guaranteed by the PIMPLE outer loop.

\begin{tcolorbox}[casebox]
  \textbf{Demonstrated by:} \caseref{case322}.
  Sub-case \texttt{sieverts\_bc} provides the $\Da \to \infty$ reference
  ($J_\text{ref}$); sub-case \texttt{surface\_recombination} runs $\Da = 1$
  and confirms $C_s = (\sqrt5-1)/2$ and the 38.2\,\% flux reduction.
\end{tcolorbox}

% ============================================================
\section{Performance Metrics}
\label{sec:metrics}
% ============================================================

\subsection{Permeation flux and speciesFlux function object}
\label{sec:flux_fo}

The area-integrated molar permeation flux through a patch is:

\begin{tcolorbox}[eqbox,title={Surface molar flux}]
\begin{equation}\label{eq:J}
  J_\text{patch}
  \;=\;
  -\!\int_\text{patch} D\,\nabla C \cdot \hat{n}\;\mathrm{d}A
  \qquad [\mol\,\s^{-1}]
\end{equation}
\end{tcolorbox}

The \texttt{speciesFlux} function object evaluates this integral at every
write time without modifying the solver.  It outputs both the area-averaged
flux density $[\mol\,\m^{-2}\,\s^{-1}]$ and the total flux $[\mol\,\s^{-1}]$
to \texttt{postProcessing/<name>/<time>/speciesFlux.dat}.

The diffusivity field \texttt{D\_<species>} is registered as
\texttt{AUTO\_WRITE} by the \texttt{speciesModel}, making it available to
the function object at every write time.

\textbf{Dictionary syntax:}
\begin{verbatim}
functions
{
    wallPermeation
    {
        type        speciesFlux;
        libs        ("libspeciesPost.so");
        region      wall;
        species     C_H2;
        patches     (wall_to_hitec);
        writeControl writeTime;
    }
}
\end{verbatim}

\begin{tcolorbox}[casebox]
  \textbf{Demonstrated by:} \caseref{case320}
  (permeation flux through SS316 wall into Hitec coolant),
  \caseref{case321}, \caseref{case322} (permeation flux at purge face).
\end{tcolorbox}

\subsection{Permeation reduction factor}
\label{sec:prf}

The permeation reduction factor (PRF) quantifies the effectiveness of a
barrier coating:

\begin{tcolorbox}[eqbox,title={Permeation reduction factor}]
\begin{equation}\label{eq:prf}
  \mathrm{PRF}
  \;=\;
  \frac{J_\text{without barrier}}{J_\text{with barrier}}
  \;=\;
  \frac{C_0/R_\text{total,bare}}{C_0/R_\text{total,coated}}
\end{equation}
\end{tcolorbox}

For a two-layer composite (coating thickness $L_c$, substrate thickness
$L_s$) with equal $K_s$ on both sides, the steady-state PRF can be
approximated by the ratio of total diffusion resistances.  High PRF arises
when the coating has a much larger $K_s/D$ ratio (lower permeability) than
the substrate.

The permeability of a material is $\Phi = K_s^2 D$ (in
$[\mol\,\m^{-1}\,\s^{-1}\,\Pa^{-1}]$, the Richardson--Dushman formula for
diatomic gases).

\begin{tcolorbox}[casebox]
  \textbf{Demonstrated by:} \caseref{case319}.
  A tungsten carbide (WC) coating on SS316 at $T = 1073\,\K$:
  \[
    \mathrm{PRF} \approx 200
  \]
  confirmed by comparing the \texttt{with\_coating} and \texttt{no\_coating}
  sub-cases.  The PRF is dominated by the large $K_s$ ratio ($K_{s,\text{WC}} \gg
  K_{s,\text{SS316}}$), which lowers the dissolved concentration in the coating
  and therefore reduces the flux into the substrate.
\end{tcolorbox}

% ============================================================
\section{Membrane Equilibrium Boundary Conditions}
\label{sec:membc}
% ============================================================

Two additional boundary conditions support membrane distillation modelling
where the species concentration and the thermal boundary condition at the
membrane face are coupled through the vapour pressure.

\subsection{Antoine equilibrium (\texttt{antoineEquilibrium})}
\label{sec:antoine}

The \bctype{antoineEquilibrium} boundary condition is a
\texttt{fixedValueFvPatchScalarField} that computes the equilibrium vapour
molar concentration from the local face temperature using the Antoine equation:

\begin{tcolorbox}[eqbox,title={Antoine vapour pressure}]
\begin{equation}\label{eq:antoine}
  \ln\bigl(p_\text{vap}\,[\Pa]\bigr)
  \;=\; A \;-\; \frac{B}{C_\text{ant} + T}
\end{equation}
\end{tcolorbox}

The face concentration is set using the ideal-gas law:

\begin{equation}\label{eq:antoine_C}
  C_\text{face} = \frac{p_\text{vap}}{R\,T} \qquad [\mol\,\m^{-3}]
\end{equation}

The temperature $T$ is read from the registered temperature field on the same
patch.  Because $T$ at the membrane face changes as the flow solution
converges, \bctype{antoineEquilibrium} is evaluated every PIMPLE outer
corrector, providing a dynamic coupling between concentration and temperature.

\textbf{Dictionary syntax:}
\begin{verbatim}
C_H2O_at_membrane
{
    type        antoineEquilibrium;
    A           23.197;    // ln(Pa) form of Antoine equation
    B           3885.7;    // K
    C           -42.98;    // K  (negative shifts reference)
    value       uniform 0;
}
\end{verbatim}

This BC replaces \texttt{codedFixedValue} used in earlier cases and is
physically identical but compiled and registered as a proper named BC.

\begin{tcolorbox}[casebox]
  \textbf{Demonstrated by:} \caseref{case325}, \caseref{case326},
  \caseref{case327} (membrane face $C_\text{H$_2$O}$, both membrane–feed and
  membrane–permeate interfaces).
\end{tcolorbox}

\subsection{Latent heat flux (\texttt{latentHeatFlux})}
\label{sec:latent}

The \bctype{latentHeatFlux} boundary condition is a
\texttt{mixedFvPatchScalarField} applied to the temperature field $T$ in
fluid channel regions at their membrane-facing patch.  It adds the latent heat
of phase change to the thermal balance:

\begin{tcolorbox}[eqbox,title={Latent heat coupling}]
\begin{equation}\label{eq:latent}
  q_\text{latent}
  \;=\;
  J \cdot L_\text{vap} \cdot M
  \qquad [\W\,\m^{-2}]
\end{equation}
\end{tcolorbox}

where $J\,[\mol\,\m^{-2}\,\s^{-1}]$ is the vapour molar flux through the
membrane (computed from the neighbour region's diffusivity and concentration
gradient), $L_\text{vap}\,[\J\,\mol^{-1}]$ is the molar latent heat, and
$M\,[\kg\,\mol^{-1}]$ is the molar mass.

The BC contributes $q_\text{latent}$ to the heat balance as a
\texttt{refGrad} correction: on the evaporation (feed) side, heat is
\emph{removed} by evaporation; on the condensation (permeate) side, heat is
\emph{added}.  The value fraction follows the same conjugate-transfer formula
as \bctype{coupledTemperature}, ensuring flux continuity for the thermal field.

\textbf{Dictionary syntax:}
\begin{verbatim}
feed_at_membrane
{
    type        latentHeatFlux;
    Lvap        44000;          // J/mol  (water at 60°C)
    M           0.018015;       // kg/mol
    CName       C_H2O;          // concentration field driving the flux
    side        evaporation;    // or: condensation
    value       uniform 333;
}
\end{verbatim}

\begin{tcolorbox}[casebox]
  \textbf{Demonstrated by:} \caseref{case326}, \caseref{case327} (replaces
  \bctype{coupledTemperature} at membrane faces in feed and permeate regions).
\end{tcolorbox}

% ============================================================
\section{Solver Selection Guide}
\label{sec:solverselect}
% ============================================================

Choosing between \texttt{speciesSolid} and \texttt{speciesFluid} for a fluid
channel region is the most important design decision in setting up a new case.
The following criteria apply.

\subsection{Use \texttt{speciesSolid} for a fluid region when}

\begin{itemize}[itemsep=4pt]

  \item \textbf{The velocity field is known in advance.}
        Plug flow ($\mathbf{U} = \text{const}$), an analytically prescribed
        Poiseuille profile, or a uniform-bulk-velocity approximation.
        Set a \texttt{fixedValue} (or \texttt{uniformFixedValue}) $\mathbf{U}$
        field in \texttt{0/<region>/U} and add the appropriate
        \texttt{div(phi,C\_<species>)} and \texttt{div(rhoPhi,e)} entries to
        \texttt{fvSchemes}.

  \item \textbf{The primary research question is species or heat transport,
        not flow development.}
        When advection is simply a forcing term and its precise spatial
        distribution is not critical, the analytical-velocity approximation
        reduces cost and complexity without meaningful loss of accuracy.

  \item \textbf{The effective diffusivity already encodes flow effects.}
        In turbulent channels, the turbulent diffusivity
        $D_\text{eff} \approx \nu_t / \mathrm{Sc}_t$ can be orders of magnitude
        larger than the molecular value.  Using a constant $D_\text{eff}$ in
        \texttt{speciesSolid} captures the bulk transport without resolving the
        velocity field (e.g.\ case323, $D_\text{eff} \approx 10^{-4}\,\m^2\,\s^{-1}$).

  \item \textbf{The region is a solid.}
        \texttt{speciesSolid} uses \texttt{solidThermo} for temperature, which
        is the correct formulation for metallic walls, membranes, and barriers.
        No momentum equation is required.

\end{itemize}

\subsection{Use \texttt{speciesFluid} for a fluid region when}

\begin{itemize}[itemsep=4pt]

  \item \textbf{The velocity profile is unknown and must be computed.}
        Developing flow (entry region), complex geometry, or pressure-driven
        flow in a channel of arbitrary cross-section where the parabolic/
        turbulent profile is not prescribed.

  \item \textbf{Concentration polarisation from a developing boundary layer
        must be captured.}
        At high Schmidt number ($\mathrm{Sc} = \nu/D \gg 1$), the concentration
        boundary layer at the permeation wall is much thinner than the momentum
        BL.  Resolving $\delta_C \sim L\,\mathrm{Pe}^{-1/3}$ requires the
        correct velocity profile from the solved PIMPLE loop.

  \item \textbf{The pressure field is physically meaningful.}
        Cases where the pressure drop, velocity profile, or flow recirculation
        influences the result require the full incompressible Navier--Stokes
        solution.

  \item \textbf{The energy equation involves convective transport.}
        \texttt{speciesSolid} uses the \texttt{solidThermo} internal energy
        form $\partial(\rho e)/\partial t + \nabla\cdot\mathbf{q}$, which is
        appropriate for solids.  For a fluid, the temperature advection term
        $\nabla\cdot(\mathbf{U}\,T)$ is essential; \texttt{speciesFluid}
        provides it via the scalar transport form~\eqref{eq:fluid_T}.

\end{itemize}

\subsection{Use \texttt{compressibleSpeciesFluid} for a fluid region when}

\begin{itemize}[itemsep=4pt]

  \item \textbf{Density varies significantly with temperature or pressure.}
        Gases with large $\Delta T$ across the channel (e.g.\ helium at
        700\,K, 1\,atm or FLiBe molten-salt at 873\,K) require a proper
        compressible energy equation where pressure-work, kinetic-energy
        flux, and variable density all couple to the velocity field.

  \item \textbf{The equation of state is non-trivial.}
        \texttt{compressibleSpeciesFluid} uses \texttt{heRhoThermo} with
        user-specified EOS (\texttt{rhoConst}, \texttt{perfectGas}, etc.),
        allowing the density--temperature--pressure coupling to be resolved
        consistently.

  \item \textbf{Sensible internal energy (rather than temperature) must be
        the primary energy variable.}
        The \texttt{fluid} base solves the full compressible energy equation
        in terms of $e$~\eqref{eq:energy_comp}, including the
        pressure-work term $\nabla\cdot(\varphi\,p/\rho)$.  The temperature
        $T$ is then recovered via the EOS; this avoids operator-splitting
        errors that arise when $T$ is used as the primary variable in a
        compressible flow.

  \item \textbf{Cross-solver validation against \texttt{speciesFluid} is
        needed.}
        When two different thermophysical frameworks (incompressible vs.\
        compressible) should give the same species flux under equivalent
        conditions, running both and comparing the permeation flux is a
        strong verification tool (see \caseref{case329}--\caseref{case330}).

\end{itemize}

\begin{tcolorbox}[notebox]
  \textbf{Rule of thumb.}  For laminar channel flow at $\mathrm{Pe} \lesssim 100$,
  advection and diffusion compete on similar length scales and the velocity
  profile matters.  Use \texttt{speciesFluid} (incompressible) or
  \texttt{compressibleSpeciesFluid} (variable-density gas or liquid).  For
  $\mathrm{Pe} \gg 100$ in engineering contexts where only the bulk flux is
  needed, a turbulent-effective-$D$ in \texttt{speciesSolid} is usually
  sufficient and much cheaper.
\end{tcolorbox}

% ============================================================
\section{Numerical Implementation}
\label{sec:num}
% ============================================================

\subsection{Finite-volume discretisation}

The species equation~\eqref{eq:species} is discretised using the
standard OpenFOAM finite-volume method on an unstructured collocated mesh:

\begin{itemize}[itemsep=2pt]
  \item \textbf{Time integration:} first-order implicit Euler
        (\texttt{ddtSchemes: Euler}).
  \item \textbf{Diffusion:} second-order Gauss linear with explicit non-orthogonal
        correction (\texttt{laplacianSchemes: Gauss linear corrected}).
  \item \textbf{Gradient:} Gauss linear (\texttt{gradSchemes: Gauss linear}).
\end{itemize}

These settings are appropriate for diffusion-dominated transport on
orthogonal structured meshes (all tutorial cases use 1-D block meshes).

\subsection{PIMPLE outer loop and Picard convergence}
\label{sec:pimple}

The nonlinear boundary conditions (\texttt{quadratic}, \texttt{sqrt},
\texttt{surfaceRecombination}) are linearised by Picard (fixed-point)
iteration within the PIMPLE outer corrector loop:

\begin{enumerate}[itemsep=2pt]
  \item \textbf{Outer iteration $k$:} the nonlinear BC evaluates the
        interface condition using cell values from iteration $k-1$ (frozen
        coefficients).
  \item \textbf{Species equation:} assembled and solved with the frozen
        coefficients; cell values are updated.
  \item \textbf{Convergence check:} if the species residual $< \varepsilon$
        (\texttt{residualControl}), the outer loop exits; otherwise another
        outer iteration is performed.
\end{enumerate}

The mixed-BC weight~\eqref{eq:weight_sqrt} uses the \emph{linearised
derivative} of the partition function evaluated at the current iterate,
ensuring that flux continuity is enforced at every Picard step and that the
method converges to the exact partition condition as $k \to \infty$.

\subsection{Coupled-patch data exchange}

At each outer PIMPLE iteration, the coupled patches exchange:
\begin{itemize}[itemsep=2pt]
  \item \textbf{Neighbour cell concentration} $C_{n,\text{cell}}$
        (used to evaluate the partition refValue),
  \item \textbf{Neighbour diffusivity} $D_n$ and
        \textbf{inverse distance} $\delta_n^{-1}$
        (used to evaluate $\text{nbrKD}$ for the weight).
\end{itemize}
Data exchange uses OpenFOAM's mapped-patch infrastructure with an incremented
MPI message tag to avoid collisions when multiple coupled fields are
exchanged in the same iteration.

\subsection{Semi-implicit trapping ODE}

The McNabb--Foster ODE~\eqref{eq:trap} is solved \emph{cell by cell} with the
per-cell semi-implicit Euler update~\eqref{eq:trap_euler}.  This is performed
explicitly before the species equation assembly, so the sink term
$\partial C_t/\partial t$ is available for the right-hand side.  No global
matrix inversion is required for the ODE update.

\subsection{Library architecture}

\begin{center}
\small
\begin{tabular}{lll}
\toprule
Library & Purpose & Key classes \\
\midrule
\texttt{libspeciesTransport.so} & Core physics & \texttt{speciesModel}, \texttt{arrheniusProperty},\\
& & \texttt{McNabbFoster}, \texttt{noTrapping} \\
\texttt{libspeciesCoupling.so} & Boundary conditions & \texttt{sievertsCoupledMixed},\\
& & \texttt{surfaceRecombination},\\
& & \texttt{antoineEquilibrium},\\
& & \texttt{latentHeatFlux} \\
\texttt{libspeciesSolid.so} & Solver: solid / prescribed-flow & \texttt{speciesSolid} \\
\texttt{libspeciesFluid.so} & Solver: incompressible fluid & \texttt{speciesFluid} \\
\texttt{libcompressibleSpeciesFluid.so} & Solver: compressible fluid & \texttt{compressibleSpeciesFluid} \\
\texttt{libspeciesPost.so} & Post-processing & \texttt{speciesFlux} \\
\bottomrule
\end{tabular}
\end{center}

All cases must list the required libraries in \texttt{system/controlDict}.
Cases using only solid regions:
\begin{verbatim}
libs ("libspeciesTransport.so" "libspeciesCoupling.so"
      "libspeciesSolid.so"    "libspeciesPost.so");
\end{verbatim}
Cases with one or more \texttt{speciesFluid} regions also add:
\begin{verbatim}
libs (...  "libspeciesFluid.so");
\end{verbatim}
Cases with one or more \texttt{compressibleSpeciesFluid} regions use instead:
\begin{verbatim}
libs (...  "libcompressibleSpeciesFluid.so");
\end{verbatim}

% ============================================================
\section{Tutorial Guide}
\label{sec:tutorials}
% ============================================================

Twenty tutorial cases cover the full range from pure-diffusion verification
to fully coupled compressible fluid-flow species transport.  Cases 311--322 use
\texttt{speciesSolid} exclusively; cases 323--326 extend \texttt{speciesSolid}
with prescribed convection; cases 327--328 introduce \texttt{speciesFluid};
cases 329--330 use \texttt{compressibleSpeciesFluid}.

Each case contains an \texttt{Allrun} script, an \texttt{Allclean} script,
a per-case \texttt{README.md}, and (for verification cases) a Python
validation script in \texttt{validation/} that compares the simulation output
to an analytical or benchmark solution (exit code 0 on pass, 1 on fail).

\subsection{Case summary table}

% Column widths chosen so total = \linewidth:
%   case col  1.8cm  (fits "case330" in bold colour)
%   solver col 5.0cm  (wraps "compressibleSpeciesFluid + speciesSolid")
%   physics col = remainder (auto-computed via \dimexpr)
%   plus 4\tabcolsep inter-column gaps = 4×6pt = 24pt ≈ 0.84cm
\small\setlength{\LTcapwidth}{\linewidth}
\begin{longtable}{%
  >{\raggedright\arraybackslash}p{1.8cm}
  >{\raggedright\arraybackslash}p{5.0cm}
  p{\dimexpr\linewidth-1.8cm-5.0cm-4\tabcolsep\relax}}
\caption{Tutorial cases and the physics they demonstrate.}
\label{tab:cases}\\
\toprule
Case & Solver(s) & Physics / Key feature \\
\midrule
\endfirsthead
\toprule
Case & Solver(s) & Physics / Key feature \\
\midrule
\endhead
\midrule\multicolumn{3}{r}{\small\itshape continued on next page\ldots}\\
\endfoot
\bottomrule
\endlastfoot

\caseref{case311} & \texttt{solid} &
  Pure 1-D diffusion, step BC.  Temperature field used as species proxy.
  Validates erfc front propagation (exact analytical solution). \\[2pt]

\caseref{case312} & \texttt{speciesSolid} &
  1-D diffusion, non-uniform IC (\texttt{setFields}).
  Fourier cosine-series analytical solution.  First use of the species equation. \\[2pt]

\caseref{case313} & \texttt{speciesSolid} &
  McNabb--Foster trapping (§\ref{sec:trap}).
  Weak ($\alpha_t/D \ll 1$) and strong ($\alpha_t/D \gg 1$) trapping regimes.
  Validates semi-implicit ODE~\eqref{eq:trap_euler}. \\[2pt]

\caseref{case314} & \texttt{speciesSolid} &
  Two-region composite membrane; Sieverts--Sieverts linear partition
  (§\ref{sec:bc_linear}), $K_{s,A}/K_{s,B} = 4$.
  Validates interface concentration jump. \\[2pt]

\caseref{case315} & \texttt{speciesSolid} &
  Isothermal membrane distillation; Henry-law solubility, constant $D$.
  Linear steady-state profile. \\[2pt]

\caseref{case316} & \texttt{speciesSolid} &
  Membrane distillation with Arrhenius $D(T)$.  Temperature gradient drives
  non-linear steady-state profile.  Full energy--species coupling. \\[2pt]

\caseref{case317} & \texttt{speciesSolid} &
  Reverse osmosis: two-region feed/membrane.
  Henry partition at feed face ($K_H = 0.5$).  Solution-diffusion model. \\[2pt]

\caseref{case318} & \texttt{speciesSolid} &
  Code-to-code benchmark vs Pasler et al.\ and Fourier sine series.
  Constant BCs; transient diffusion; exact analytical solution. \\[2pt]

\caseref{case319} & \texttt{speciesSolid} &
  WC coating + SS316; real Arrhenius $D(T)$ for both materials.
  PRF $\approx 200$ (§\ref{sec:prf}).  Two sub-cases with/without coating. \\[2pt]

\caseref{case320} & \texttt{speciesSolid} &
  Three-region shell-tube HX: FLiBe $|$ SS316 $|$ Hitec.
  Conjugate heat + H$_2$ permeation; Arrhenius $D(T)$.
  \texttt{speciesFlux} FO measures $J$. \\[2pt]

\caseref{case321} & \texttt{speciesSolid} &
  Henry--Sieverts partition (§\ref{sec:bc_nonlinear}).
  Two sub-cases: \texttt{sieverts\_sieverts} and \texttt{henry\_sieverts}.
  Validates 23.6\,\% flux increase; golden-ratio interface~\eqref{eq:hs_ss}. \\[2pt]

\caseref{case322} & \texttt{speciesSolid} &
  Surface recombination / dissociation (§\ref{sec:bc_recomb}).
  Da$=1$: $C_s = (\sqrt{5}-1)/2$; 38.2\,\% flux reduction vs Sieverts BC. \\[2pt]

\caseref{case323} & \texttt{speciesSolid} &
  2-D counter-flow FLiBe--SS316--Hitec HX.
  Prescribed turbulent plug flow in fluid channels via registered $\mathbf{U}$ field.
  Turbulent-effective $D_\text{eff}$ absorbs flow effects. \\[2pt]

\caseref{case324} & \texttt{speciesSolid} &
  DCMD validation vs Khalifa 2017.  Single-region PTFE membrane;
  membrane face $C$ set from Antoine equation as \texttt{fixedValue}.
  $J_\text{sim} = 31.1$ vs $J_\text{exp} = 35\,\text{L\,m}^{-2}\,\text{h}^{-1}$. \\[2pt]

\caseref{case325} & \texttt{speciesSolid} &
  DCMD with CFD temperature polarisation.  Three-region 2-D; prescribed plug
  flow in feed/permeate; \bctype{antoineEquilibrium} at membrane faces
  (§\ref{sec:antoine}).  Temperature polarisation coefficient (TPC) emerges
  from coupled CFD rather than a heat-transfer correlation. \\[2pt]

\caseref{case326} & \texttt{speciesSolid} &
  DCMD with latent heat coupling.
  \bctype{latentHeatFlux} BC (§\ref{sec:latent}) replaces
  \bctype{coupledTemperature} at membrane faces; adds $q = J\,L_\text{vap}\,M$
  to the thermal balance.  Full membrane-distillation energy accounting. \\[2pt]

\caseref{case327} & \texttt{speciesFluid} + \texttt{speciesSolid} &
  DCMD with solved incompressible flow.  Feed/permeate: \texttt{speciesFluid}
  (PIMPLE $\mathbf{U}/p$ + scalar $T$ + $C$).  Membrane: \texttt{speciesSolid}.
  Demonstrates the complete \texttt{speciesFluid} workflow with
  \bctype{antoineEquilibrium} and \bctype{latentHeatFlux} at coupled faces. \\[2pt]

\caseref{case328} & \texttt{speciesFluid} &
  Solver validation: single-region 2-D channel, water at Re$=20$.
  Three quantitative checks: Poiseuille profile (L$_2 < 2\,\%$),
  pressure gradient ($<5\,\%$), Graetz Nu$_\infty = 7.5407$ ($<5\,\%$). \\[2pt]

\caseref{case329} & \texttt{compressibleSpeciesFluid} + \texttt{speciesSolid} &
  FLiBe molten-salt channel (873\,K) with distributed H$_2$ source $\to$
  SS316 permeation.  \texttt{heRhoThermo}/\texttt{rhoConst} EOS; compressible
  PIMPLE; $\mathrm{Pe} = 1000$; Sieverts linear interface; vacuum outer BC.
  Validates equivalence of compressible and incompressible formulations for
  constant-density liquids. \\[2pt]

\caseref{case330} & \texttt{compressibleSpeciesFluid} + \texttt{speciesSolid} &
  He gas channel (700\,K, 1\,atm) with distributed H$_2$ source $\to$
  SS316 permeation.  Identical $D$, $K_s$, source, and inlet velocity to
  \caseref{case329}.  Cross-solver validation: confirms that the compressible
  solver gives the same permeation flux as the rhoConst-liquid case across
  a 10$\times$ change in density. \\

\end{longtable}

\subsection{Physics-to-tutorial cross-reference}

\begin{table}[H]
\centering
\caption{Which tutorials exercise which physical models (all 20 cases).
  $\checkmark$ = feature present; blank = not applicable.}
\label{tab:cross}
\scriptsize
\setlength{\tabcolsep}{4pt}
\renewcommand{\arraystretch}{1.3}
\resizebox{\linewidth}{!}{%
\begin{tabular}{lrrrrrrrrrrrrrrrrrrrr}
\toprule
Physical model &
\rotatebox{75}{\caseref{311}} & \rotatebox{75}{\caseref{312}} &
\rotatebox{75}{\caseref{313}} & \rotatebox{75}{\caseref{314}} &
\rotatebox{75}{\caseref{315}} & \rotatebox{75}{\caseref{316}} &
\rotatebox{75}{\caseref{317}} & \rotatebox{75}{\caseref{318}} &
\rotatebox{75}{\caseref{319}} & \rotatebox{75}{\caseref{320}} &
\rotatebox{75}{\caseref{321}} & \rotatebox{75}{\caseref{322}} &
\rotatebox{75}{\caseref{323}} & \rotatebox{75}{\caseref{324}} &
\rotatebox{75}{\caseref{325}} & \rotatebox{75}{\caseref{326}} &
\rotatebox{75}{\caseref{327}} & \rotatebox{75}{\caseref{328}} &
\rotatebox{75}{\caseref{329}} & \rotatebox{75}{\caseref{330}} \\
\midrule
Species diffusion eq.~\eqref{eq:species}
  &\checkmark&\checkmark&\checkmark&\checkmark&\checkmark&\checkmark&\checkmark&\checkmark&\checkmark&\checkmark&\checkmark&\checkmark
  &\checkmark&\checkmark&\checkmark&\checkmark&\checkmark&\checkmark&\checkmark&\checkmark \\
Energy eq.~\eqref{eq:energy}
  &\checkmark&\checkmark&\checkmark&\checkmark&\checkmark&\checkmark&\checkmark&\checkmark&\checkmark&\checkmark&\checkmark&\checkmark
  &\checkmark&\checkmark&\checkmark&\checkmark&\checkmark&\checkmark&\checkmark&\checkmark \\
Arrhenius $D(T)$~\eqref{eq:arrhenius}
  & & & & & &\checkmark& & &\checkmark&\checkmark& &
  &\checkmark&\checkmark&\checkmark&\checkmark&\checkmark& & & \\
Sieverts law~\eqref{eq:sieverts}
  & & &\checkmark&\checkmark& & & &\checkmark&\checkmark&\checkmark&\checkmark&\checkmark
  &\checkmark& & & & & &\checkmark&\checkmark \\
Henry law~\eqref{eq:henry}
  & & & & &\checkmark&\checkmark&\checkmark& & & &\checkmark&
  & &\checkmark&\checkmark&\checkmark&\checkmark& & & \\
McNabb--Foster~\eqref{eq:trap}
  & & &\checkmark& & & & & & & & &
  & & & & & & & & \\
Linear partition~\eqref{eq:linear_part}
  & & & &\checkmark& & & &\checkmark&\checkmark&\checkmark&\checkmark&
  &\checkmark& & & & & &\checkmark&\checkmark \\
Sqrt partition~\eqref{eq:sqrt_part}
  & & & & & & & & & & &\checkmark&
  & & & & & & & & \\
Quadratic partition~\eqref{eq:quad_part}
  & & & & & & &\checkmark& & & &\checkmark&
  & & & & & & & & \\
Surface recombination~\eqref{eq:recomb}
  & & & & & & & & & & & &\checkmark
  & & & & & & & & \\
PRF metric~\eqref{eq:prf}
  & & & & & & & & &\checkmark& & &
  & & & & & & & & \\
speciesFlux FO~\eqref{eq:J}
  & & & & & & & & & &\checkmark&\checkmark&\checkmark
  &\checkmark&\checkmark&\checkmark&\checkmark&\checkmark& & & \\
Multi-region ($\geq 2$ regions)
  & & & & & & & & & &\checkmark& &
  &\checkmark& &\checkmark&\checkmark&\checkmark& &\checkmark&\checkmark \\
Antoine equilibrium BC (§\ref{sec:antoine})
  & & & & & & & & & & & &
  & &\checkmark&\checkmark&\checkmark&\checkmark& & & \\
Latent heat flux BC (§\ref{sec:latent})
  & & & & & & & & & & & &
  & & & &\checkmark&\checkmark& & & \\
Incompressible PIMPLE~\eqref{eq:mom}
  & & & & & & & & & & & &
  & & & & &\checkmark&\checkmark& & \\
Compressible PIMPLE~\eqref{eq:mom_comp}
  & & & & & & & & & & & &
  & & & & & & &\checkmark&\checkmark \\
Volumetric source $S_\text{vol}$~\eqref{eq:species}
  & & & & & & & & & & & &
  & & & & & & &\checkmark&\checkmark \\
\bottomrule
\end{tabular}}
\end{table}

\subsection{Running the tutorials}

\textbf{Prerequisites.} Source OpenFOAM-13 and build the libraries:
\begin{verbatim}
source /opt/openfoam13/etc/bashrc
cd $WM_PROJECT_USER_DIR/multiSpeciesRegionFoam
./Allwmake 2>&1 | tee build.log
\end{verbatim}

\textbf{Individual case:}
\begin{verbatim}
cd tutorials/case314-composite-membrane
./Allrun
python3 validation/membrane_compare.py --plot
\end{verbatim}

\textbf{Validation exit code:} each validation script exits 0 on
\textsc{pass} (flux within 5\,\% of analytical) and 1 on \textsc{fail}.


% ============================================================
\newpage
\section{Notation Summary}
\label{sec:notation}
% ============================================================

\begin{center}
\begin{tabular}{lll}
\toprule
Symbol & Meaning & SI units \\
\midrule
$C$               & Mobile species concentration       & $\mol\,\m^{-3}$ \\
$C_t$             & Trapped concentration               & $\mol\,\m^{-3}$ \\
$D(T)$            & Diffusivity (Arrhenius)              & $\m^2\,\s^{-1}$ \\
$e$               & Sensible internal energy (\texttt{compressibleSpeciesFluid}) & $\J\,\kg^{-1}$ \\
$E_a$             & Activation energy                   & $\J\,\mol^{-1}$ \\
$f_t$             & Trap fraction                        & --- \\
$J$               & Molar flux (area-averaged)           & $\mol\,\m^{-2}\,\s^{-1}$ \\
$K = \tfrac{1}{2}|\mathbf{U}|^2$ & Specific kinetic energy (\texttt{compressibleSpeciesFluid}) & $\J\,\kg^{-1}$ \\
$K_d(T)$          & Dissociation rate                   & $\mol\,\m^{-2}\,\s^{-1}\,\Pa^{-1}$ \\
$K_H(T)$          & Henry solubility constant            & $\mol\,\m^{-3}\,\Pa^{-1}$ \\
$K_r(T)$          & Recombination rate                  & $\m^4\,\mol^{-1}\,\s^{-1}$ \\
$K_s(T)$          & Sieverts solubility constant         & $\mol\,\m^{-3}\,\Pa^{-1/2}$ \\
$L$               & Layer thickness                      & $\m$ \\
$N$               & Molar trap-site density              & $\mol\,\m^{-3}$ \\
$n_t = f_t N$     & Available trap density               & $\mol\,\m^{-3}$ \\
$p$               & Pressure (thermodynamic in compressible, kinematic in incompressible) & $\Pa$ or $\m^2\,\s^{-2}$ \\
$R$               & Gas constant ($8.314$)               & $\J\,\mol^{-1}\,\K^{-1}$ \\
$S_\text{vol}$    & Volumetric species source            & $\mol\,\m^{-3}\,\s^{-1}$ \\
$T$               & Temperature                          & $\K$ \\
$\mathbf{U}$      & Fluid velocity                       & $\m\,\s^{-1}$ \\
$w$               & Mixed-BC value fraction              & --- \\
$X_0$             & Arrhenius pre-exponential factor     & (same as $X$) \\
$\alpha_d(T)$     & Detrapping rate                     & $\s^{-1}$ \\
$\alpha_t$        & Trapping rate                       & $\s^{-1}$ \\
$\delta^{-1}$     & Cell-centre to face distance (inv.) & $\m^{-1}$ \\
$\varepsilon$     & Trap binding energy                  & eV \\
$\kappa(T)$       & Thermal conductivity                 & $\W\,\m^{-1}\,\K^{-1}$ \\
$\rho$            & Fluid density                        & $\kg\,\m^{-3}$ \\
$\varphi$         & Face mass flux ($= \rho\mathbf{U}\cdot\mathbf{S}_f$, \texttt{phi} in OpenFOAM) & $\kg\,\s^{-1}$ \\
$\varphi_\text{vol}$ & Face volumetric flux ($= \varphi/\rho_f$, \texttt{phiU}) & $\m^3\,\s^{-1}$ \\
$\Phi = K_s^2 D$  & Permeability (diatomic gas)          & $\mol\,\m^{-1}\,\s^{-1}\,\Pa^{-1}$ \\
$\Da$             & Damköhler number (surface kinetics) & --- \\
$\mathrm{PRF}$    & Permeation reduction factor          & --- \\
\bottomrule
\end{tabular}
\end{center}


% ============================================================
\addcontentsline{toc}{section}{References}

\begin{thebibliography}{9}

\bibitem{Hattab2024}
  Hattab, N., Siriano, S., Giannetti, F.
  \textit{An OpenFOAM multi-region solver for tritium transport modeling
  in fusion systems.}
  \textit{Fusion Engineering and Design} \textbf{202} (2024) 114362.

\bibitem{McNabb1963}
  McNabb, A., Foster, P.\,K.
  \textit{A new analysis of the diffusion of hydrogen in iron and ferritic steels.}
  \textit{Trans.\ Metall.\ Soc.\ AIME} \textbf{227} (1963) 618--627.

\bibitem{Pick1985}
  Pick, M.\,A., Sonnenberg, K.
  \textit{A model for atomic hydrogen-metal interactions --- application to
  recycling, recombination and permeation.}
  \textit{J.\ Nucl.\ Mater.} \textbf{131} (1985) 208--220.

\bibitem{Anderl1992}
  Anderl, R.\,A.\ et al.
  \textit{Hydrogen transport behavior of metal coatings for plasma-facing components.}
  \textit{J.\ Nucl.\ Mater.} \textbf{196--198} (1992) 986--991.

\bibitem{Baker2012}
  Baker, R.\,W.
  \textit{Membrane Technology and Applications}, 3rd~ed.
  Wiley, 2012.

\bibitem{Ockwig2007}
  Ockwig, N.\,W., Nenoff, T.\,M.
  \textit{Membranes for hydrogen separation.}
  \textit{Chem.\ Rev.} \textbf{107} (2007) 4078--4110.

\bibitem{CFDDirect}
  CFD Direct.
  \textit{Modular Solvers in OpenFOAM.}
  \url{https://cfd.direct/openfoam/free-software/modular-solvers/}

\end{thebibliography}

\end{document}
